\subsection{线性微分方程积分的存在性}

设 E 是域 R 上的完备赋范空间，J 是 R 中不退化为一点的区间。称微分方程

\[
{\frac {d \mathbf {x}}{d t}} = \mathbf {f} (t, \mathbf {x})\tag{1}
\]

其中 \(\mathbf{f}\) 定义在 \(\mathbf{J} \times \mathbf{E}\) 上，称为线性方程，如果对每个 \(t \in \mathbf{J}\)，映射 \(\mathbf{x} \mapsto \mathbf{f}(t, \mathbf{x})\) 是从 E 到其自身中的连续仿射线性映射 \(^1\) ；若令 \(\mathbf{b}(t) = \mathbf{f}(t, 0)\)，则映射 \(\mathbf{x} \mapsto \mathbf{f}(t, \mathbf{x}) - \mathbf{f}(t, 0) = \mathbf{f}(t, \mathbf{x}) - \mathbf{b}(t)\) 便是从 E 到其自身中的连续线性映射；往后我们把这一映射记作 \(A(t)\)，并写 \(A(t).\mathbf{x}\)（或简写 \(A(t)\mathbf{x}\) ）表示它在点 \(\mathbf{x} \in E\) 的值；这样，线性微分方程 (1) 就可写成

\[
{\frac {d \mathbf {x}}{d t}} = A (t). \mathbf {x} + \mathbf {b} (t)\tag{2}
\]

其中 \(\mathbf{b}\) 是从 J 到 E 中的映射；当 \(\mathbf{b} = 0\) 时，称线性微分方程 (2) 为齐次的。

例. 1) 当 E 在 R 上是有限维 n 时，可以把自同态 \(A(t)\) 与它关于 E 的任意基的矩阵 \(\left(a_{ij}(t)\right)\) 等同起来（《代数》，II，第 343 页）；当把向量 \(x \in E\) 与它关于所考虑的 E 的基的分量的列矩阵 \(\left(x_{j}\right)\) 等同时，表达式 \(A(t)\) .x 符合《代数》中的一般约定（《代数》，II，第 343 页，命题 2）。在这种情形，方程 (2) 等价于数量微分方程组

\[
\frac {d x _ {i}}{d t} = \sum_ {j = 1} ^ {n} a _ {i j} (t) x _ {j} + b _ {i} (t) \quad (1 \leqslant i \leqslant n).\tag{3}
\]

\begin{enumerate}
\def\labelenumi{\arabic{enumi})}
\setcounter{enumi}{1}
\tightlist
\item
  设 \(\mathbf{G}\) 是 \(\mathbf{R}\) 上的完备赋范代数，而 \(\mathbf{a}(t)\) 、\(\mathbf{b}(t)\) 与 \(\mathbf{c}(t)\) 是从 \(\mathbf{J}\) 到 \(\mathbf{G}\) 中的三个映射；方程
\end{enumerate}

\[
\frac {d \mathbf {x}}{d t} = \mathbf {a} (t) \mathbf {x} + \mathbf {x b} (t) + \mathbf {c} (t)
\]

是线性微分方程；这里 \(A(t)\) 是线性映射 \(\mathbf{x} \mapsto \mathbf{a}(t)\mathbf{x} + \mathbf{x}\mathbf{b}(t)\) （从 \(G\) 到其自身）。

对每个 \(t \in J\)，\(A(t)\) 是由 E 到其自身的连续线性映射（E 的连续自同态）所成之集 \(\mathcal{L}(\mathrm{E})\) 的一个元素；我们知道（《一般拓扑学》，X，第 298 页），\(\mathcal{L}(\mathrm{E})\) 赋以范数 \(\|U\| = \sup_{\|x\| \leqslant 1} \|Ux\|\) 后是域 R 上的完备赋范代数，并且 \(\|UV\| \leqslant \|U\|\) \(\|V\|\) 。

在本节中我们始终假定下列条件满足：

\begin{enumerate}
\def\labelenumi{\alph{enumi})}
\item
  映射 \(t \mapsto A(t)\)（从 \(\mathbf{J}\) 到 \(\mathcal{L}(\mathrm{E})\) 中）是规则的。
\item
  映射 \(t \mapsto \mathbf{b}(t)\)（从 J 到 E 中）是规则的。
\end{enumerate}

当 \(\mathrm{E}\) 的维数为 \(n\) 时，\(\mathcal{L}(\mathrm{E})\) （作为拓扑向量空间）同构于 \(\mathbf{R}^{n^2}\)，而条件 \(a\) ）表示 \(a_{ij}(t)\)（矩阵 \(A(t)\) 的元素）都是 \(\mathrm{J}\) 上的规则函数。

由于 \(\left\| A(t')\mathbf{x} - A(t)\mathbf{x}\right\| \leqslant \left\| A(t') - A(t)\right\| \| \mathbf{x}\|\) ，映射

\[
t \mapsto A (t). \mathbf {x} + \mathbf {b} (t)
\]

对每个 \(\mathbf{x} \in \mathrm{E}\) 都是规则的；此外，

\[
\left\| A (t) \mathbf {x} _ {1} - A (t) \mathbf {x} _ {2} \right\| = \left\| A (t) \left(\mathbf {x} _ {1} - \mathbf {x} _ {2}\right) \right\| \leqslant \left\| A (t) \right\| \left\| \mathbf {x} _ {1} - \mathbf {x} _ {2} \right\|
\]

对任何 \(t \in J\) 与 E 中的 \(x_{1}\) 、\(x_{2}\) 都成立；换言之，(2) 的右端满足 IV 第 165 页引理 1 的条件，并且关于规则函数 \(\|A(t)\|\) 在 \(J \times E\) 上是利普希茨的。由此得到（IV，第 173 页，推论 2）：

定理 1. 设 \(t \mapsto A(t)\) 是从 \(\mathbf{J}\) 到 \(\mathcal{L}(\mathrm{E})\) 中的规则映射，并设 \(t \mapsto \mathbf{b}(t)\) 是从 \(\mathbf{J}\) 到 \(\mathrm{E}\) 中的规则映射。对 \((t_0, \mathbf{x}_0)\) 的每个点 \(\mathbf{J} \times \mathbf{E}\)，线性方程 (2) 有唯一一个定义在整个 \(\mathbf{J}\) 上且等于 \(\mathbf{x}_0\) 于点 \(t_0\) 的解。

